
\chapter{Glossar}

\begin{description}
  \item[\IX{Arteriosklerose}] Systemerkrankung der
  Arterien, welche zu Ablagerung von Fetten, Bindegewebe und Kalk in
  den Blutgefäßwänden und somit zur Schädigung und
  Veränderungen der Gefäßen führt.
%  \item[\IX{bergen}]
  \item[\IX{Cave}]   \Lang{lat.} \Speech{\enquote{Hüte Dich vor
  \dots}}. Auch: \enquote{Achtung, gib Acht!}
  \item[\IX{Computertomographie}] \Lang{Syn.} \I{CT}.
  Röntgenuntersuchung, bei welcher Schnittbilder des
  Körpers erstellt werden. Dabei rotiert  eine
  Röntgenröhre um den Patienten und sendet
  fächerförmig Röntgenstrahlen durch den Körper. Auf der
  gegenüberliegenden Seite  befindet sich ein Detektor,
  welcher die eintreffenden Röntgenstrahlen registriert. Aus
  den Messwerten berechnet ein Computer
  Querschnittsbilder der untersuchten Region. Moderne
  Systeme können Bilder in verschiedenen Ebenen
  Rekonstruieren, bzw. 3D-Darstellungen erzeugen. Bei der
  Untersuchung von Weichteilstrukturen wird häufig ein
  Kontrastmittel (\FullPrettyRef{topic:kontrasmittel})
  eingesetzt.
  \item[\IX{Erysipel}] \Lang{Sny.} Rotlauf, Wundrose. \label{topic:erysipel}
  Infektionserkrankung, v.\,a. der unteren Extremität, welche
  gekennzeichnet ist durch eine lokale Rötung und Schwellung 
  andere betreffende Extremität. Ursächlich ist das
  Eindringen von Keimen, zumeist Streptokokken oder
  Staphylokokken, durch eine Eintrittspforte (zum Beispiel
  Insektenstiche, Fußpilz, \ldots). Ein Rotlauf kam mit 
  Fieber einhergehen. Eine antibiotische Behandlung ist
   erforderlich, in schweren Fällen muss diese
  antibiotische Behandlung stationär mittels intravenöse
  Infusion erfolgen.
  \item[\IX{Gestose}] Oberbegriff für durch die
  Schwangerschaft bedingte Erkrankungen. Man unterscheidet
  zwischen Gestosen der Frühschwangerschaft (im
  1. Trimenon, \I{Frühgestosen}) und der Spätschwangerschaft
  (3. Trimenon, \I{Spätgestosen}). Darunter fallen die
  \I{Hyperemesis gravidarum} (übermäßiges, anhaltendes
  Erbrechen), die EPH-Gestose bzw. Präekplampsie
  (\ref{topic:eph-gestose}), die Eklampsie
  (\ref{topic:eklampsie}) und das HELLP-Syndrom
  (\ref{topic:hellp-Syndrom}).
  \item[\IX{Herzkatheteruntersuchung}]
  \label{topic:Herzkatheteruntersuchung} Interventionelle
  röntgengestützte Untersuchung des Herzens bzw. der
  Herzkrankgefäße. Dabei wird über eine Arterie
  (klassisch über die Leistenarterie) ein Katheter
  eingebracht und bis zum Herz bzw. bis knapp vor das Herz
  zu den Abgängen der Koronargefäße aus der Aorta vorgeschoben.
  Anschließend wird ein Kontrastmittel appliziert um den
  Ventrikel (\I{Ventrikulographie}) bzw. die Gefäße
  (\I{Koronarangiographie}) mittels Röntgen sichtbar zu
  machen.
  
  Bei der Koronarangiographie werden die Herzkranzgefäße auf
  signifikante Engstellen, welche das Substrat einer
  koronaren Herzkrankheit, eines akuten Koronarsyndroms bzw.
  eines Herzinfarktes darstellen, untersucht. Im Rahmen des
  Vorganges können diese Engstellen auch gleich behandelt
  werden (\II{PCI}, \I{Percutaneous Coronary Intervention};
  \II{PTCA}, \I{Perkutaneous Transluminal
  Coronary Angioplasty}).
  Als Optionen stehen die Dehnung mittels eines Ballons und/oder die Einbringung eines \E{Stents}
  in das betroffene Gefäß zur Wahl. Ein \T{Stent} ist vom
  Prinzip aufgebaut wie ein sich selbst aufspannendes,
  schlauchförmiges Metallgitter, welches in das Gefäß
  eingebracht wird und dieses an der betreffenden Stelle
  aufdehnt. Es gibt unbeschichtete (\I{BMS}, \I{Bare Metal
  Stent}) und beschichtete (\I{DES}, \I{Drug Eluting
  Stent}) Stents.
  
  \begin{PictureSeriesWide}{}{Bilderserie:
    \EEE{Koronargefäße}}
    \PictureSeriesRowWide{3}
    {./images/hirtler-aass/Herzvorderflaeche.pdf}
    {Das Herz mit seinen Koronargefäßen \SourcePicture{Hirtler}}
    {./images/wikipedia-cc-by/Hk_coronary_bionerd.jpg}
    {Darstellung der
    Herzkranzgefäße während einer Herzkatheteruntersuchung.
    \SourcePicture{WmCo
    \protect\enquote{Bionerd}, CC-BY-3.0}}
    {./images/wikipedia-cc-by/Heart_coronary_artery_lesion-lq.jpg}
    {Die Koronargefäße versorgen den
    Herzmuskel von außen nach innen. \SourcePicture{Patrick J.
    Lynch, CC-BY-2.5}} {./images/wikipedia-cc-by/Hk_coronary_bionerd.jpg}
    {empty}
    {}
  \end{PictureSeriesWide}
  \item[\IX{Kontrastmittel}] \label{topic:kontrasmittel}
  Kontrastmittel verbessern die Darstellung von Strukturen
  bei bildgebenden Verfahren (Röntgendiagnostik inkl.
  Computertomographie (CT) und Angiographie,
  Magnetresonanztomografie (MRT) und Sonografie (Ultraschall). Kontrastmittel können
  unerwünschte Nebenwirkung haben. Darunter fallen
  allergische (anaphylaktische) Reaktionen,
  Schilddrüsenüberfunktion (bei Jod-haltigen
  Kontrastmitteln), sowie eine Einschränkung der
  Nierenfunktion. Häufige unbedenkliche Nebenwirkungen sind
  ein Wärme- oder Hitzegefühl sowie ein metallischer
  Geschmack.
  \item[\IX{Tokolyse}] \Lang{Syn.} \I{Wehenhemmung}. Hemmung
  der Wehentätigkeit durch Medikamente. Eine Tokolyse wird
  vor allem bei vorzeitiger Wehentätigkeit durchgeführt,
  wenn eine Geburt aufgrund der Unreife des Kindes oder
  anderen Gründen verzögert werden soll. Ebenso kann eine
  Tokolyse durchgeführt werden um eine bevorstehende Geburt
  soweit zu verzögernd, bis die Patientin in eine geeignete
  Einrichtung verlegt werden kann.
%  \item[\IX{retten}]
  \item[\IX{Rotlauf}] siehe Erysipel, \FullPrettyRef{topic:erysipel}
  \item[\IX{Ruptur}] Riss, Einreissen, Durchbruch.
%  \item[\IX{Sofortbehandlung}]
%  \item[\IX{Sofortmaßnahmen, lebensrettende}]
  \item[Stent] Siehe Herzkatheteruntersuchung
  (\FullPrettyRef{topic:Herzkatheteruntersuchung})
  \item[\IX{Sukkurs}] Hilfe, Unterstützung; Einheit, die
  als Verstärkung bzw. Unterstützung eingesetzt ist.
  \item[\IX{Wehenhemmung}] siehe Tokolyse
\end{description}
